% Seating the microcontroller. This is identical for every project, so the
% coordinates live here rather than in any one project chapter. The diagram
% generator reads the \bbhole definitions below for every breadboard diagram
% it draws, which is why this file is the only place they appear.
\subsubsection{Attach the microcontroller to the breadboard}
Orient the microcontroller so that the pin labeled \textbf{5V} goes into hole \bbhole{mcu.5v}{H1}
(Row H, Column 1), \textbf{GPIO2} into \bbhole{mcu.gpio2}{D1}, \textbf{GPIO20} into
\bbhole{mcu.gpio20}{H7}, and \textbf{GPIO21} into \bbhole{mcu.gpio21}{D7}. Then press it down---this
might take more force than you expect.

\begin{figure}[H]
    \centering
    \maybeincludegraphics[width=.95\linewidth]{common/microcontroller_seated_in_breadboard.pdf}
    \caption{So far, so good!}
\end{figure}

Once it is seated, each microcontroller pin shares a terminal strip with the free holes in its
column. For example, \textbf{GPIO 20} sits in \textbf{H7}, so a jumper plugged into \textbf{I7}
or \textbf{J7} connects to GPIO 20. That is why every wiring table in this guide starts its
jumpers in columns 1--7: rows I and J for the top pins, and rows A--C for the bottom pins. The
board's body covers the rest of the holes in those columns.

\begin{center}
\begin{tabular}{|c|l|l|}
    \hline
    \textbf{Column} & \textbf{Top pin (row H)} & \textbf{Bottom pin (row D)} \\
    \hline
    1 & 5V      & GPIO 2  \\
    2 & GND     & GPIO 3  \\
    3 & 3V3     & GPIO 4  \\
    4 & GPIO 10 & GPIO 5  \\
    5 & GPIO 9  & GPIO 6  \\
    6 & GPIO 8  & GPIO 7  \\
    7 & GPIO 20 & GPIO 21 \\
    \hline
\end{tabular}
\end{center}
